\documentclass{article}
\usepackage{geometry}
\geometry{
	a4paper,
	left=2.5cm,   
	right=2.5cm,   
	top=2.5cm,      
	bottom=2.5cm,  
	headheight=13.6pt
}

% 最简单、最稳定的中文方案（自动适配系统，永不报错）
\usepackage[UTF8, scheme=plain]{ctex}

% 设置英文默认字体为 Times New Roman
\usepackage{fontspec}
\setmainfont{Times New Roman}

\usepackage{amsmath}
\usepackage{amssymb}
\usepackage{booktabs}
\usepackage{caption}
\usepackage{setspace}%设置间距
\onehalfspacing
\setlength{\parskip}{4pt}

\usepackage{array}
\usepackage{listings}
\usepackage{graphicx}
\usepackage{enumitem}
\usepackage{subcaption}
\usepackage{multirow}
\usepackage{mathrsfs}
\usepackage{bm}
\usepackage{xcolor}
\usepackage{colortbl}
\usepackage{nccmath}
\everymath{\displaystyle}  % 全文行内公式自动启用 displaystyle

% 定义代码颜色
\definecolor{codegreen}{rgb}{0,0.6,0}
\definecolor{NavyBlue}{RGB}{0,0,128}
\definecolor{PineGreen}{RGB}{1,121,111}

% 配置 listings 环境
\lstset{
	language=Python,
	tabsize=4,
	captionpos=b,
	numbers=left,
	numbersep=1em,
	sensitive=true,
	showtabs=false,
	frame=shadowbox,
	breaklines=true,
	keepspaces=true,
	showspaces=false,
	showstringspaces=false,
	breakatwhitespace=false,
	basicstyle=\ttfamily\small,
	keywordstyle=\color{NavyBlue},
	commentstyle=\color{codegreen},
	numberstyle=\color{black},
	stringstyle=\color{PineGreen!90!black},
	rulesepcolor=\color{red!20!green!20!blue!20}
}

\definecolor{LightGreen}{rgb}{0.9, 0.95, 0.9}
\definecolor{LightBlue}{RGB}{229, 238, 255}
\definecolor{LightPink}{RGB}{255, 229, 237}

% 修正表格标题格式：居中对齐，字体为黑体
\captionsetup[table]{
	justification=centering,
	labelsep=quad,
	font={bf}
}

% 设置图片标题（图注）字体为黑体
\captionsetup[figure]{
	justification=centering,
	labelsep=quad,
	font={bf}
}

% 设置子图标题字体为黑体
\captionsetup[subfigure]{
	font={bf}
}

\begin{document}
	
	\begin{center}
		{\LARGE \textbf{概率论期末复习}} \\[0.5cm]
		{\large {\kaishu 石继楷 \quad 编写}} \\[0.3cm]
		{\large 2026年6月7日}
	\end{center}
	
	\section{第一章\quad 概率空间与随机事件}
	\subsection{$\bm{Def}$ 超几何分布（不放回）}
	\[
	P(X=k)=\frac{\mathrm{C}_r^k \mathrm{C}_b^{n-k}}{\mathrm{C}_{r+b}^n}
	\]
	\subsection{$\bm{Def}$ 二项分布（放回）}
	\[
	P(X=k)=\frac{\mathrm{C}_n^k r^k b^{n-k}}{(r+b)^n}
	\]
	\subsection{$\bm{Def}$ 可测空间 $(\Omega,\mathcal{F},P)$}
	$\mathcal{F}$ 为$\sigma$代数(事件域)，$P$为概率测度。
	\subsection{$\bm{Thm}$ 概率的五条性质}
	\begin{enumerate}
		\item 可列可加：$P\Big(\bigcup_{k=1}^{\infty}A_k\Big)=\sum_{k=1}^{\infty}P(A_k),\ A_k$ 两两不交。
		\item 加法公式：$P\Big(\bigcup_{k=1}^nA_k\Big)=\sum P(A_k)-\sum_{i<j}P(A_i\cap A_j)+\sum_{i<j<k}P(A_iA_jA_k)+\cdots+(-1)^{n-1}P(A_1A_2\cdots A_n)$
		\item 下连续性：$P\Big(\bigcup_{k=1}^{\infty}A_k\Big)=\lim_{k\to\infty}P(A_k),\ A_n\subset A_{n+1}$
		\item 上连续性：$P\Big(\bigcap_{k=1}^{\infty}A_k\Big)=\lim_{k\to\infty}P(A_k),\ A_n\supset A_{n+1}$
		\item 次可加：$P\Big(\bigcup_{k=1}^{\infty}A_k\Big)\le\sum_{k=1}^{\infty}P(A_k),\ \forall A_n\in\mathcal{F}$
	\end{enumerate}
	\subsection{$\bm{Def}$ 测度$(\Omega,\mathcal{F})$上概率测度$\mu$}
	\begin{enumerate}
		\item 非负：$\forall A\in\mathcal{F},\mu(A)\ge0$
		\item 空集测度为$0$：$\mu(\emptyset)=0$
		\item 可列可加：$\mu\Big(\bigcup_{k=1}^{\infty}A_k\Big)=\sum_{k=1}^{\infty}\mu(A_k)$
		\item 有限：$\mu(\Omega)<+\infty$
	\end{enumerate}
	\subsection{$\bm{Thm}$ 全概率公式}
	\[
	P(A)=\sum_{n=1}^{\infty}P(B_n)P(A\mid B_n)
	\]
	\subsection{$\bm{Thm}$ 多个事件独立}
	$\forall 2\le k\le n,\forall 1\le i_1<i_2<\cdots<i_k\le n$
	\[
	P(A_{i_1}A_{i_2}\cdots A_{i_k})=P(A_{i_1})P(A_{i_2})\cdots P(A_{i_k})
	\]
	共$2^n-n-1$个式子。
	
	\section{第二章\quad 随机变量及其分布}
	\subsection{$\bm{Def}$ 随机变量}
	$\xi:\Omega\to\mathbb{R},\forall B\in\mathcal{B}(\mathbb{R}),\xi^{-1}(B)\in\mathcal{F}$
	\subsection{$\bm{Def}$ 分布函数 $G(x)=\mu(-\infty,x]$}
	\begin{enumerate}
		\item 右连续：$\lim\limits_{x\to x_0^+}F(x)=F(x_0)$
		\item 单调不降
		\item 有界$\iff \mu$有限测度
		\item 规范性：$\lim\limits_{x\to-\infty}F(x)=0,\lim\limits_{x\to+\infty}F(x)=1$
	\end{enumerate}
	\[
	P(\xi\in(a,b])=F(b)-F(a)
	\]
	\subsection{$\bm{Def}$ 二项分布$B(n,p)$}
	\[
	E\xi=np,\quad P(\xi=k)=\mathrm{C}_n^kp^k(1-p)^{n-k}
	\]
	最可能值$k=[(n+1)p]$或$k-1$。
	\subsection{$\bm{Def}$ 几何分布$G(p)$}
	\[
	E\xi=\sum_{i=1}^{\infty}iP(\xi=i)=\frac1p,\quad P(\xi=k)=(1-p)^{k-1}p
	\]
	无记忆性：$P(\xi>k+n\mid\xi>n)=P(\xi>k)$
	\subsection{$\bm{Def}$ 帕斯卡分布$Pa(n,p)$，$\xi_n$第$n$次成功试验次数}
	\[
	P(\xi_n=k)=\mathrm{C}_{k-1}^{n-1}p^n(1-p)^{k-n}
	\]
	\subsection{$\bm{Def}$ 泊松分布$P(\lambda)$}
	\[
	E\xi=\lambda,\quad P(\xi=k)=\frac{\lambda^k}{k!}e^{-\lambda}
	\]
	\subsection{$\bm{Def}$ 泊松过程$\{N_t,t\ge0\}$}
	\begin{enumerate}
		\item 独立增量：$\{N_{t_1}-N_{t_0}=k_1\},\{N_{t_3}-N_{t_2}=k_2\}$独立
		\item 平稳增量：$P(N_{s+t}-N_s=k)=P(N_t=k)$
		\item 普通性：$1-P(N_t=0)-P(N_t=1)=o(t)$
	\end{enumerate}
	\[
	P_n(t)=\frac{(\lambda t)^n}{n!}e^{-\lambda t}
	\]
	\subsection{$\bm{Def}$ 均匀分布$U(a,b)$}
	\[
	p(x)=
	\begin{cases}
		\displaystyle\frac1{b-a},&a<x<b\\[4pt]
		0,&\text{其他}
	\end{cases},\quad
	F(x)=
	\begin{cases}
		0,&x<a\\[4pt]
		\displaystyle\frac{x-a}{b-a},&a\le x\le b\\[4pt]
		1,&x>a
	\end{cases}
	\]
	\subsection{$\bm{Def}$ 指数分布$E(\lambda)$}
	\[
	E\xi=\frac1\lambda,\quad p(x)=\lambda e^{-\lambda x}(x>0),\quad F(x)=1-e^{-\lambda x}(x>0)
	\]
	无记忆性。
	\subsection{$\bm{Def}$ $\Gamma$分布$\Gamma(\alpha,\lambda)$}
	\[
	E\xi=\frac\alpha\lambda,\quad
	p(x)=\frac{\lambda^\alpha x^{\alpha-1}}{\Gamma(\alpha)}e^{-\lambda x}=\frac{\lambda^k}{(k-1)!}x^{k-1}e^{-\lambda x}
	\]
	\[
	F(x)=\sum_{k=n}^{\infty}\frac{(\lambda t)^k}{k!}e^{-\lambda t},\quad \Gamma(s)=\int_{0}^{+\infty}x^{s-1}e^{-t}dt
	\]
	\subsection{$\bm{Def}$ 正态分布$N(a,\sigma^2)$}
	\[
	p(x)=\frac1{\sqrt{2\pi}\sigma}\exp\left\{-\frac{(x-a)^2}{2\sigma^2}\right\},\quad F(x)=\int_{-\infty}^xp(y)dy=\Phi\left(\frac{x-a}{\sigma}\right)
	\]
	$\Phi(\cdot)$为标准正态$N(0,1)$。
	\subsection{$\bm{Def}$ 不降函数的广义逆$F^{-1}(y)=\inf\{x\mid F(x)\ge y\}$}
	\begin{enumerate}
		\item $\forall y\in(0,1),F(F^{-1}(y))\ge y$
		\item $F^{-1}$在$(0,1)$不降、左连续
		\item $F$在$x$跳跃，则$F^{-1}$在$[F(x-),F(x)]$为常数
		\item $F$在$(a,b)$严格增，则$F^{-1}$在$y$不取小于$b-a$的跳跃
	\end{enumerate}
	
	\section{第三章\quad 随机向量与独立性}
	\subsection{$\bm{Def}$ 联合分布函数}
	\[
	F(x,y)=P(\xi\le x,\eta\le y)
	\]
	性质：矩形增量$\Delta F(a,b)=F(b_1,b_2)-F(b_1,a_2)-F(a_1,b_2)+F(a_1,a_2)\ge0$。
	\subsection{$\bm{Def}$ 联合密度}
	\begin{enumerate}
		\item 非负：$\forall x,y,p(x,y)\ge0$
		\item 归一：$\displaystyle \iint_{\mathbb{R}^2}p(x,y)dxdy=1$
	\end{enumerate}
	\subsection{$\bm{Def}$ 边缘分布}
	离散：
	\[
	p_{\xi}(x_i)=\sum_jp(x_i,y_j),\quad p_{\eta}(y_j)=\int_{-\infty}^{+\infty}p(x,y)dx
	\]
	连续：
	\[
	F_{\xi}(x)=\sum_{x_i\le x}p_{\xi}(x_i),\quad F_{\eta}(y)=\int_{-\infty}^yp_{\eta}(v)dv
	\]
	\subsection{$\bm{Def}$ 条件密度}
	\[
	p_{\eta\mid\xi}(y\mid x)=\frac{p(x,y)}{p_{\xi}(x)}
	\]
	\subsection{$\bm{Thm}$ 独立判定}
	离散：$p(x_i,y_j)=p_{\xi}(x_i)p_{\eta}(y_j)$
	
	\noindent 连续：$p(x,y)=p_{\xi}(x)p_{\eta}(y)$
	\subsection{$\bm{Thm}$ 卷积公式}
	离散：
	\[
	P(\xi+\eta=n)=\sum_{k=0}^nP(\xi=k)P(\eta=n-k)
	\]
	连续：
	\[
	p_{\xi+\eta}(z)=\int_{-\infty}^{+\infty}p(x,z-x)dx
	\]
	\subsection{$\bm{Def}$ 期望}
	离散：$\displaystyle E\xi=\sum_{x_i<x} x_ip(x_i)$
	
	\noindent 连续：$E\displaystyle \xi=\int_{-\infty}^{+\infty}xp(x)dx$
	\[
	E(\xi^2-\xi)=\int_{-\infty}^{+\infty}(x^2-x)p(x)dx
	\]
	\subsection{$\bm{Thm}$ 期望性质}
	$\xi,\eta$独立$\implies E(\xi\eta)=E\xi E\eta$；
	不加独立：$E(\xi+\eta)=E\xi+E\eta,\ E(X-C)=EX-C$。
	\subsection{$\bm{Thm}$ 积分收敛定理}
	\begin{enumerate}
		\item 有界收敛：若$\xi_n\xrightarrow{\text{a.s.}}\xi,|\xi_n|\le M$，则$\lim\limits_{n\to\infty} E\xi_n=E(\lim\limits_{n\to\infty}\xi_n)$
		\item Levi单调收敛：$\{\xi_n\}$单调不降$\xi_n\to\xi\ \text{a.s.}\implies\lim\limits_{n\to\infty} E\xi_n=E(\lim\limits_{n\to\infty}\xi_n)$
		\item Fatou引理：$E(\varliminf\limits_{n\to\infty}\xi_n)\le\varliminf\limits_{n\to\infty} E\xi_n$
		\item 控制收敛：$\exists\eta$可积，$|\xi_n|\le\eta\implies E(\lim\limits_{n\to\infty}\xi_n)=E\xi$
	\end{enumerate}
	
	\section{第四章\quad 方差、协方差、条件期望}
	\subsection{$\bm{Thm}$ Cauchy-Schwarz不等式}
	\[
	\big(E(\xi\eta)\big)^2\le E\xi^2\cdot E\eta^2
	\]
	等号$\iff \eta=t\xi\ \text{a.s.}$
	\subsection{$\bm{Def}$ 方差}
	\[
	\mathrm{Var}\xi=E(\xi-E\xi)^2=E\xi^2-(E\xi)^2
	\]
	\begin{enumerate}
		\item $\mathrm{Var}\xi\ge0$，等号$\iff\xi=E\xi\ \text{a.s.}$
		\item $\mathrm{Var}(c\xi)=c^2\mathrm{Var}\xi$
		\item $\mathrm{Var}(\xi+c)=\mathrm{Var}\xi$
		\item $\xi,\eta$独立$\implies\mathrm{Var}(\xi+\eta)=\mathrm{Var}\xi+\mathrm{Var}\eta$
		\item $f(c)=E(\xi-c)^2$，$c=E\xi$取最小
	\end{enumerate}
	\subsection{$\bm{Def}$ 协方差}
	\[
	\mathrm{Cov}(\xi,\eta)=E(\xi\eta)-E\xi E\eta
	\]
	\begin{enumerate}
		\item 对称：$\mathrm{Cov}(\xi,\eta)=\mathrm{Cov}(\eta,\xi)$
		\item 双线性：$\mathrm{Cov}(a\xi+b\eta,\zeta)=a\mathrm{Cov}(\xi,\zeta)+b\mathrm{Cov}(\eta,\zeta)$
		\item $\mathrm{Var}(\xi+\eta)=\mathrm{Var}\xi+\mathrm{Var}\eta+2\mathrm{Cov}(\xi,\eta)$
	\end{enumerate}
	\subsection{$\bm{Def}$ 相关系数}
	\[
	r=\frac{\mathrm{Cov}(\xi,\eta)}{\sqrt{\mathrm{Var}\xi}\sqrt{\mathrm{Var}\eta}}
	\]
	\subsection{$\bm{Def}$ 原点矩$m_k=E\xi^k$，中心矩$C_k=E(\xi-E\xi)^k$}
	\subsection{$\bm{Def}$ 条件期望}
	离散：
	\[
	E(f(\xi,\eta)\mid\eta=y_j)=\sum_{i=1}^{\infty}f(x_i,y_j)\frac{p(x_i,y_j)}{p_{\eta}(y_j)}
	\]
	连续：
	\begin{align}
	\displaystyle E(f(\xi,\eta)\mid\eta=y)&=\int_{-\infty}^{+\infty}f(x,y)\frac{p(x,y)}{p_{\eta}(y)}dx\\
	&=\int_{-\infty}^{+\infty}f(x,y)\frac{p(x,y)}{\int_{-\infty}^{+\infty} p(x,y)dx}dx
	\end{align}
	\subsection{$\bm{Def}$ 条件方差}
	\[
	\mathrm{Var}(\xi\mid\eta=y)=E(\xi^2\mid\eta=y)-\big(E(\xi\mid\eta=y)\big)^2
	\]
	\subsection{$\bm{Thm}$ 全期望公式}
	\[
	E(f(\xi,\eta))=\sum_nP(B_n)E(f(\xi,\eta)\mid B_n)
	\]
	平滑性：$E\big(E(f(\xi,\eta)\mid\eta=y)\big)=E(f(\xi,\eta))$
	Wald定理：$E\Big(\sum_{k=1}^{\tau}\xi_k\Big)=E\tau\cdot E\xi_1$
	
	\section{第五章\quad 特征函数}
	\subsection{$\bm{Def}$ 特征函数$f(t)=E(e^{it\xi})$}
	离散：$\displaystyle f(t)=\sum_{k=1}^{\infty}e^{itx_k}P(\xi=x_k)$
	
	\noindent 连续：$\displaystyle f(t)=\int_{-\infty}^{+\infty}e^{itx}p(x)dx$
	\subsection{$\bm{Thm}$ 矩与特征函数}
	\[
	E\xi^k=i^{-k}f^{(k)}(0),\quad E\xi=i^{-1}f'(0),\quad E\xi^2=-f''(0)
	\]
	\subsection{$\bm{Prop}$ 特征函数性质}
	\begin{enumerate}
		\item $|f(t)|\le f(0)=1$
		\item 共轭对称：$f(-t)=\overline{f(t)}$
		\item $f(t)$在$\mathbb{R}$一致连续
		\item 线性变换：$\eta=a+b\xi\implies f_{\eta}(t)=e^{ita}f_{\xi}(bt)$
	\end{enumerate}
	唯一性定理：分布由特征函数唯一确定。
	
	\section{第六章\quad 随机变量四种收敛}
	\subsection{$\bm{Def}$ 点点收敛}
	$\forall\omega,\xi_n(\omega)\to\xi(\omega)\iff\forall\varepsilon>0,\exists N,n>N\implies|\xi_n(\omega)-\xi(\omega)|<\varepsilon$
	\subsection{$\bm{Def}$ 依概率收敛}
	\[
	\lim_{n\to\infty}P\big(|\xi_n-\xi|\ge\varepsilon\big)=0,\quad \text{记}\ \xi_n\stackrel{P}{\to}\xi
	\]
	\subsection{$\bm{Def}$ 几乎处处收敛}
	\[
	P\big(\omega:\lim_{n\to\infty}\xi_n(\omega)=\xi(\omega)\big)=1
	\iff \lim_{k\to\infty}P\Big(\bigcup_{n=k}^{\infty}\big\{|\xi_n-\xi|\ge\varepsilon\big\}\Big)=0
	\]
	记$\xi_n\xrightarrow{\text{a.s.}}\xi$
	\subsection{$\bm{Def}$ $L^r$收敛}
	$E|\xi|^r<+\infty\implies E|\xi_n|^r<+\infty,\lim\limits_{n\to\infty}E|\xi_n-\xi|^r=0$
	\subsection{$\bm{Def}$ 依分布收敛}
	$C_F$为$F$连续点，$\forall x\in C_F,\lim\limits_{n\to\infty}F_n(x)=F(x)$，记$F_n\stackrel{d}{\to}F$；
	
	\noindent 几乎处处收敛$\implies$依概率收敛$\implies$依分布收敛。
	
	\section{第七章\quad 大数定律与中心极限定理}
	\subsection{$\bm{Def}$ 弱大数定律}
	\[
	\frac1n\sum_{k=1}^n(\xi_k-E\xi_k)\stackrel{P}{\to}0,\quad n\to\infty
	\]
	\subsection{$\bm{Def}$ 强大数定律}
	\[
	\frac1n\sum_{k=1}^n(\xi_k-E\xi_k)\xrightarrow{\text{a.s.}}0
	\]
	
	\subsection{$\bm{Thm}$ Markov不等式}
	\[P(|\xi|\geq \varepsilon)\leq \dfrac{E|\xi|^r}{\varepsilon^r}\]
	
	\subsection{$\bm{Thm}$ 切比雪夫不等式}
	\[
	P(|\xi-E\xi|\ge\varepsilon)\le\frac1{\varepsilon^2}\mathrm{Var}\xi
	\]
	\subsection{$\bm{Thm}$ 切比雪夫大数律}
	$\{\xi_n\}$两两不相关、方差有界，则满足弱大数。
	\subsection{$\bm{Thm}$ Bernoulli大数律}
	$X_n$为$n$次伯努利成功次数，$\displaystyle\frac{X_n}{n}\stackrel{P}{\to}p$。
	\subsection{$\bm{Lemma}$ Borel-Cantelli引理}
	\begin{enumerate}
		\item $\sum P(A_n)<\infty\implies P(\varlimsup\limits_{n\to\infty} A_n)=0$
		\item $\{A_n\}$独立，$\sum P(A_n)=\infty\implies P(\varlimsup\limits_{n\to\infty} A_n)=1$
	\end{enumerate}
	\subsection{$\bm{Thm}$ Borel强大数律}
	$\displaystyle\frac{X_n}{n}\xrightarrow{\text{a.s.}}p$
	\subsection{$\bm{Thm}$ 科尔莫戈罗夫强大数律}
	$\{\xi_n\}$独立同分布，$E|\xi_1|<+\infty$，则$\displaystyle\frac1n\sum_{k=1}^n(\xi_k-E\xi_k)\xrightarrow{\text{a.s.}}0$
	\subsection{$\bm{Def}$ 中心极限定理}
	$\{\xi_n\}$独立，$E\xi_n=a,\mathrm{Var}\xi_n=\sigma^2>0$
	\[
	P(\tilde{S}_n\le x)\to\Phi(x),\quad \tilde{S}_n=\frac{\sum_{k=1}^n(\xi_k-E\xi_k)}{\sqrt{\mathrm{Var}\Big(\sum_{k=1}^n\xi_k\Big)}},\ E\tilde{S}_n=0,\mathrm{Var}\tilde{S}_n=1
	\]
	\subsection{$\bm{Thm}$ Lindeberg-Levy定理}
	$\{\xi_n\}$独立同分布，$E\xi_1=a,0<\mathrm{Var}\xi_1=\sigma^2<+\infty$，满足中心极限定理。
	
\end{document}